% TOP_PROBLEM: {"id":30007646,"problem_number":"LOCAL-30007646","title":"Dynastic persistence beyond parents","category_id":21,"category":{"id":21,"name":"economics","display_name":"Economics","description":"Open economic theory statements, published research agendas, and proposed scoped specifications; question provenance and historical priority are qualified per record.","slug":"economics","order_index":21,"created_at":"2026-10-08T00:00:00Z"},"status":"research_question","external_url":null,"legacy_ids":[],"published":false,"statement":"In the explicitly specified setting: How far does economic advantage persist across three or more generations, and which direct ancestral influences survive after parental characteristics are adequately measured? The mathematical statement above fixes the target, assumptions and scope of this catalogue formulation.","economics":{"original_id":135,"field":"Inequality and economic mobility","kind":"identification","formalization_basis":"adapted_research_question","source_status":"explicit_open","priority_status":"earliest_found","priority_note":"Earliest directly inspected qualifying passage in this audit; earlier origins were not established. A review or research-agenda statement does not imply original priority.","earliest_verified_reference":{"authors":["Gary Solon"],"title":"What Do We Know So Far about Multigenerational Mobility?","year":2017,"url":"https://docs.iza.org/dp10623.pdf","doi":"10.1111/ecoj.12495","locator":"Section 2, final paragraph, printed p. 11","evidence_summary":"Future research must determine which of several observationally similar persistence mechanisms matter in which circumstances."},"current_reference":{"authors":["Gary Solon"],"title":"What Do We Know So Far about Multigenerational Mobility?","year":2017,"url":"https://docs.iza.org/dp10623.pdf","doi":"10.1111/ecoj.12495","locator":"Section 2, final paragraph, printed p. 11","evidence_summary":"Future research must determine which of several observationally similar persistence mechanisms matter in which circumstances."},"formalization_note":"The cited passage establishes a research direction, not this exact mathematical statement. The notation, population, policy menu, assumptions, and estimands are an original adaptation for this catalogue; no universal unsolved theorem or source priority is claimed."},"statusQualification":"Published question or research agenda verified at the cited locator. The mathematical specification is adapted for this catalogue and includes additional assumptions; source publication does not establish first-ever priority or certify this exact model remains unsolved. The cited passage establishes a research direction, not this exact mathematical statement. The notation, population, policy menu, assumptions, and estimands are an original adaptation for this catalogue; no universal unsolved theorem or source priority is claimed.","statusReviewedAt":"2026-10-08"}
% FIRST_POSTED: 2026-10-08
% ENTRY_KIND: statement_only
% !TeX program = lualatex
\documentclass[11pt]{article}
\usepackage{fontspec}
\usepackage[a4paper,margin=25mm]{geometry}
\usepackage{amsmath,amssymb,mathtools}
\usepackage{xurl}
\usepackage[hidelinks]{hyperref}
\newcommand{\sref}[2]{\hyperlink{source:#1:#2}{[S#2]}}
\setlength{\emergencystretch}{3em}
\begin{document}
% problemId:30007646
\section*{Dynastic persistence beyond parents}\label{30007646}
\subsection{Definitions and mathematical statement}
Consider linked families observed for at least three generations. Let \(Y_{ig}\) be a consistently measured lifetime economic-status variable for family line \(i\) in generation \(g\), and let \(\widetilde Y_{ig}=Y_{ig}+\epsilon_{ig}\) denote the observed proxy with explicitly modeled measurement error \(\epsilon_{ig}\). The descriptive task is to estimate the joint transition law \(P(Y_{ig}\mid Y_{i,g-1},Y_{i,g-2},\ldots,X_i)\), where \(X_i\) includes cohort and institutional context. Test whether more distant ancestry adds predictive information after reliably measured parental status. A nonzero grandparental regression coefficient is not itself a direct causal effect. Compare models containing direct grandparental investments, transmitted unobserved traits, group persistence, and correlated measurement error that may yield similar associations. Seek additional data or exogenous variations that distinguish their quantitative contributions. Define any causal intervention separately, for example a feasible grandparental transfer holding specified parental resources fixed. Estimate decay of ancestry associations and transportability across family structures and historical periods. No first-order autoregression or universal persistence rate is imposed. Existing linked-data findings resolve some descriptive comparisons; the remaining research agenda concerns identifying underlying mechanisms under declared measurement and selection assumptions.

\par\textbf{Formulation and scope.} The cited passage establishes a research direction, not this exact mathematical statement. The notation, population, policy menu, assumptions, and estimands are an original adaptation for this catalogue; no universal unsolved theorem or source priority is claimed.

\par\textbf{Open-question provenance.} An explicit question or research agenda was verified at the source locator. This does not by itself certify that every later version remains unresolved \sref{30007646}{1}.

\par\textbf{Historical priority.} Earliest directly inspected qualifying passage in this audit; earlier origins were not established. A review or research-agenda statement does not imply original priority.

\subsection{Short English statement}
In the explicitly specified setting: How far does economic advantage persist across three or more generations, and which direct ancestral influences survive after parental characteristics are adequately measured? The mathematical statement above fixes the target, assumptions and scope of this catalogue formulation.

\subsection{Sources}
\begin{itemize}\raggedright
\item[\textnormal{[S1]}]\hypertarget{source:30007646:1}{} Gary Solon. 2017. What Do We Know So Far about Multigenerational Mobility?. Section 2, final paragraph, printed p. 11.
\url{https://docs.iza.org/dp10623.pdf}.
DOI: 10.1111/ecoj.12495.
{\small\textit{Earliest verified question/agenda reference; historical priority qualified below. Future research must determine which of several observationally similar persistence mechanisms matter in which circumstances.}}
\end{itemize}
\end{document}
